% ============================================================
% IA BIAIS INITIATIVE — Manifeste fondateur
% Bilingue FR/EN · Design premium
% Version corrigée — Compilation garantie
% Auteur : Mohamadou SIDIBE
% Version : 1.0 · Octobre 2026
% Compilation : pdflatex (2 passes)
% ============================================================

\documentclass[11pt,a4paper,openany]{book}

\usepackage[utf8]{inputenc}
\usepackage[T1]{fontenc}
\usepackage[french]{babel}
\usepackage{lmodern}
\usepackage{microtype}
\usepackage[a4paper,margin=2cm,top=2cm,bottom=2cm]{geometry}
\usepackage{textcomp}
\usepackage{xcolor}
\usepackage{graphicx}
\usepackage{tikz}
\usetikzlibrary{shapes,positioning,calc,backgrounds,fit,decorations.pathreplacing}
\usepackage{booktabs}
\usepackage{tabularx}
\usepackage{array}
\usepackage{colortbl}
\usepackage{tcolorbox}
\tcbuselibrary{skins,breakable}
\usepackage{enumitem}
\usepackage{titlesec}
\usepackage{fancyhdr}
\usepackage{lastpage}
\usepackage{fontawesome5}
\usepackage[hidelinks]{hyperref}

\DeclareUnicodeCharacter{2122}{\texttrademark}

% ------------------------------------------------------------
% PALETTE
% ------------------------------------------------------------
\definecolor{navy}{HTML}{0F172A}
\definecolor{purple}{HTML}{4C1D95}
\definecolor{cyan}{HTML}{0891B2}
\definecolor{orange}{HTML}{D97706}
\definecolor{green}{HTML}{059669}
\definecolor{red}{HTML}{B91C1C}
\definecolor{gray}{HTML}{94A3B8}
\definecolor{light}{HTML}{F8FAFC}
\definecolor{mint}{HTML}{ECFDF5}

% ------------------------------------------------------------
% STYLE CHAPITRES
% ------------------------------------------------------------
\titleformat{\chapter}[display]
  {\normalfont\huge\bfseries\color{navy}}
  {\filright\textsc{\Large\color{purple}Chapitre \thechapter}}
  {10pt}
  {\Huge\filright}
  [\vspace{0.3em}\textcolor{purple}{\rule{\textwidth}{2pt}}]

\titlespacing*{\chapter}{0pt}{-20pt}{25pt}

\titleformat{\section}
  {\normalfont\Large\bfseries\color{navy}}
  {\thesection}{0.8em}{}
  [\vspace{-0.3em}\textcolor{cyan}{\rule{\textwidth}{0.8pt}}]

\titleformat{\subsection}
  {\normalfont\large\bfseries\color{purple}}
  {\thesubsection}{0.6em}{}

% ------------------------------------------------------------
% EN-TÊTE / PIED
% ------------------------------------------------------------
\pagestyle{fancy}
\fancyhf{}
\fancyhead[LE,RO]{\small\textcolor{navy}{\textbf{IA BIAIS INITIATIVE — Manifeste}}}
\fancyhead[RE,LO]{\small\textcolor{gray}{\nouppercase{\rightmark}}}
\fancyfoot[C]{\small\textcolor{gray}{\thepage}}
\fancyfoot[LE,RO]{\small\textcolor{gray}{contact@ia-biais-initiative.org}}
\fancyfoot[RE,LO]{\small\textcolor{gray}{v1.0 · Oct. 2026}}
\renewcommand{\headrulewidth}{0.4pt}
\renewcommand{\footrulewidth}{0.4pt}

% ------------------------------------------------------------
% BOÎTES
% ------------------------------------------------------------
\newtcolorbox{ideabox}{
  colback=mint, colframe=green, boxrule=1pt, arc=4pt,
  left=10pt,right=10pt,top=8pt,bottom=8pt, breakable,
  title={\faLightbulb\ Idée clé}, coltitle=white, colbacktitle=green,
  fonttitle=\bfseries
}

\newtcolorbox{quoteauthor}{
  colback=light, colframe=purple, boxrule=0pt, leftrule=4pt,
  arc=0pt, left=10pt,right=10pt,top=8pt,bottom=8pt, breakable
}

\newtcolorbox{keybox}{
  colback=yellow!8, colframe=orange, boxrule=1pt, arc=4pt,
  left=10pt,right=10pt,top=8pt,bottom=8pt, breakable
}

\newtcolorbox{ctabox}{
  colback=navy, colframe=navy, boxrule=0pt, arc=6pt,
  left=12pt,right=12pt,top=12pt,bottom=12pt, breakable,
  coltext=white,
}

% ------------------------------------------------------------
% MÉTADONNÉES PDF
% ------------------------------------------------------------
\hypersetup{
  pdftitle={IA BIAIS INITIATIVE — Manifeste fondateur},
  pdfauthor={Mohamadou SIDIBE},
  pdfsubject={Observatoire mondial des biais algorithmiques},
  pdfkeywords={biais algorithmiques, équité, IA responsable, gouvernance IA},
}

% ============================================================
\begin{document}

% ------------------------------------------------------------
% PAGE DE TITRE
% ------------------------------------------------------------
\begin{titlepage}
\centering
\vspace*{1cm}

\begin{tikzpicture}
  \node[fill=navy,text=white,minimum width=15cm,minimum height=4.5cm,
        align=center,inner sep=18pt,rounded corners=8pt]
    {\fontsize{34}{40}\selectfont\textbf{IA BIAIS}\\[6pt]
     \fontsize{34}{40}\selectfont\textbf{INITIATIVE}\\[10pt]
     \fontsize{14}{18}\selectfont\textcolor{gray}{Observatoire Mondial des Biais}\\[2pt]
     \fontsize{14}{18}\selectfont\textcolor{gray}{dans l'Intelligence Artificielle}};
\end{tikzpicture}

\vspace{1cm}

{\LARGE\textbf{\textcolor{purple}{MANIFESTE FONDATEUR}}\\[4pt]
 \large\textcolor{gray}{\textit{Founding Manifesto}}}

\vspace{1.2cm}

% Filet spectral
\begin{tikzpicture}
  \fill[purple] (0,0) rectangle (3,0.4);
  \fill[cyan]   (3,0) rectangle (6,0.4);
  \fill[orange] (6,0) rectangle (9,0.4);
  \fill[red]    (9,0) rectangle (12,0.4);
  \fill[green]  (12,0) rectangle (15,0.4);
\end{tikzpicture}

\vspace{1cm}

\begin{quote}
\centering
{\Large\itshape\textcolor{purple}{« L'IA sans équité n'est qu'un outil de discrimination automatisée. »}}\\[8pt]
{\large — SIDIBE Mohamadou, Fondateur}
\end{quote}

\vfill

\begin{tcolorbox}[colback=light,colframe=purple,boxrule=1pt,arc=6pt,width=13cm,center]
\centering
{\Large\textbf{Mohamadou SIDIBE}}\\[6pt]
{\normalsize Expert en Gouvernance des Données \& IA}\\[2pt]
{\normalsize Chief Data Officer · Polytechnique ExED 2026}\\[8pt]
{\normalsize\textbf{Fondateur, IA BIAIS INITIATIVE}}\\[2pt]
{\small\texttt{contact@ia-biais-initiative.org}}\\[4pt]
{\small\href{https://ia-biais-initiative.org}{ia-biais-initiative.org}}
\end{tcolorbox}

\vspace{1cm}

{\small\textcolor{gray}{
  Version 1.0 · Octobre 2026\\
  Licence CC BY 4.0 · Réutilisation libre avec attribution
}}

\end{titlepage}

% ------------------------------------------------------------
% TABLE DES MATIÈRES
% ------------------------------------------------------------
\frontmatter
\tableofcontents
\clearpage

% ------------------------------------------------------------
% PRÉAMBULE
% ------------------------------------------------------------
\chapter*{Préambule}
\addcontentsline{toc}{chapter}{Préambule}

Nous, fondateurs de l'IA BIAIS INITIATIVE, réunis autour d'une conviction partagée — que \textbf{l'équité algorithmique n'est pas un paramètre technique mais un choix de société} — déclarons ce qui suit.

\vspace{0.3cm}

\begin{ideabox}
L'intelligence artificielle est en train de devenir l'infrastructure invisible de nos sociétés. Elle décide, trie, classe, recommande, refuse, oriente. Mais elle le fait souvent \emph{sans rendre de comptes}. C'est ce vide que nous venons combler.
\end{ideabox}

\vspace{0.3cm}

Ce Manifeste expose :
\begin{enumerate}[leftmargin=1.5em]
  \item Pourquoi nous nous levons aujourd'hui
  \item Ce que disent les leaders mondiaux de l'IA
  \item Notre mission, notre vision, nos valeurs
  \item Notre réponse : la Base\_IBI
  \item Notre connexion au framework AEXIA\textsuperscript{TM}
  \item Notre gouvernance
  \item Notre feuille de route
  \item Comment nous rejoindre
\end{enumerate}

\vspace{0.5cm}

\begin{quoteauthor}
\textbf{\textcolor{purple}{Note aux lecteurs}}\\[4pt]
Ce Manifeste s'adresse à quatre publics complémentaires :
\begin{itemize}[leftmargin=1.5em, topsep=2pt, itemsep=1pt]
  \item \textbf{Chercheurs et académiques} — pour la rigueur méthodologique
  \item \textbf{Entreprises et praticiens} (CDO, DPO, DSI, Data Scientists) — pour les outils opérationnels
  \item \textbf{ONG et société civile} — pour l'ancrage citoyen et la défense des droits
  \item \textbf{Régulateurs et institutions} — pour l'alignement réglementaire
\end{itemize}
\end{quoteauthor}

\clearpage

% ============================================================
% PARTIE I
% ============================================================
\mainmatter
\part{Pourquoi nous nous levons}
\label{part:pourquoi}

\chapter{Une révolution sans garde-fous}

\section{L'ampleur du phénomène}

L'intelligence artificielle transforme chaque secteur de l'économie mondiale en une décennie : santé, justice, recrutement, crédit, sécurité, éducation, mobilité, médias. Cette transformation est porteuse de promesses immenses. Elle est aussi porteuse d'un \textbf{risque systémique invisible} : les biais algorithmiques.

Les décisions automatisées déterminent désormais :
\begin{itemize}[leftmargin=1.5em]
  \item qui obtient un crédit ou un logement
  \item qui est convoqué à un entretien d'embauche
  \item qui reçoit un diagnostic prioritaire
  \item qui est placé en détention provisoire
  \item qui est signalé à la police ou aux services sociaux
  \item qui voit son dossier médical traité en urgence
\end{itemize}

\begin{keybox}
\textbf{Ces décisions ne sont pas neutres.} Elles reproduisent, amplifient et automatisent des discriminations historiques.
\end{keybox}

\section{Ce que nous savons, et ce que nous ignorons}

Nous savons que les biais algorithmiques existent. Nous savons qu'ils affectent des millions de personnes. Nous savons qu'ils sont documentés dans tous les secteurs critiques.

Nous ne savons pas — et personne ne sait aujourd'hui — \textbf{combien de systèmes d'IA déployés dans le monde présentent des biais significatifs}. Nous n'avons pas de référentiel mondial, pas de taxonomie universellement reconnue, pas d'instrument de mesure comparable.

\begin{ideabox}
Cette lacune n'est pas un accident. Elle est un \textbf{choix} — celui des acteurs qui bénéficient du déploiement massif des systèmes d'IA sans avoir à documenter leurs biais.
\end{ideabox}

\chapter{Ce que disent les leaders mondiaux de l'IA}

Nous ne sommes pas seuls à alerter. Les voix les plus autorisées du domaine convergent.

\section{Geoffrey Hinton — Prix Nobel de Physique 2024}

\begin{quoteauthor}
\textit{« Les entreprises technologiques poursuivent des profits à court terme au détriment de la sécurité à long terme de l'humanité. Nous avons besoin d'une attention urgente et énergique pour contrer les risques de l'IA. »}\\[4pt]
\textbf{— Discours Nobel, Stockholm, 2024}
\end{quoteauthor}

Hinton, co-inventeur des réseaux de neurones profonds, a \textbf{quitté Google en 2023} pour alerter librement. Son geste est sans précédent dans l'histoire de l'informatique.

\section{Yoshua Bengio — Prix Turing 2018}

\begin{quoteauthor}
\textit{« Le public doit avoir une voix beaucoup plus forte dans les décisions qui façonneront notre avenir collectif. »}\\[4pt]
\textbf{— Appel international pour des garde-fous IA, 2024}
\end{quoteauthor}

Bengio, directeur de l'Institut Mila (Montréal), préside le Comité International d'Experts sur l'IA issu du Sommet de Bletchley Park.

\section{Sam Altman — CEO OpenAI}

\begin{quoteauthor}
\textit{« L'extinction par l'IA est un risque inacceptable. Nous ne pouvons pas nous permettre de continuer à accélérer sans avoir résolu les questions fondamentales d'alignement et de sécurité. »}\\[4pt]
\textbf{— Déclaration publique, septembre 2026}
\end{quoteauthor}

\section{Dario Amodei — CEO Anthropic}

\begin{quoteauthor}
\textit{« Accélérer sans aligner est la définition même de l'imprudence. »}\\[4pt]
\textbf{— ABC News, 2025}
\end{quoteauthor}

\section{Elon Musk — Fondateur xAI}

\begin{quoteauthor}
\textit{« L'IAG représente un risque significativement supérieur à celui des armes nucléaires. »}\\[4pt]
\textbf{— Benzinga, septembre 2026}
\end{quoteauthor}

\section{Les institutions internationales}

\begin{center}
\renewcommand{\arraystretch}{1.3}
\small
\begin{tabularx}{\textwidth}{@{}>{\bfseries}p{4cm}p{2cm}X@{}}
\toprule
\rowcolor{navy}
\textcolor{white}{Institution / Événement} & \textcolor{white}{Date} & \textcolor{white}{Message clé} \\
\midrule
\textbf{Sommet IA Paris} & Fév. 2025 & 130+ États signataires, appel à la lutte contre les biais \\
\rowcolor{light}
\textbf{Bletchley Park Summit} & Nov. 2023 & 28 États + UE, déclaration sur les risques \\
\textbf{Lettre FLI} & Mars 2023 & 1 800+ signataires, moratoire 6 mois \\
\rowcolor{light}
\textbf{Résolution ONU A/78/L.49} & Mars 2024 & Première résolution ONU sur la gouvernance IA \\
\textbf{EU AI Act} & Août 2024 & Sanctions jusqu'à 35 M€ / 7 \% du CA mondial \\
\rowcolor{light}
\textbf{RGPD art. 22} & En vigueur & Droit à l'explication algorithmique \\
\bottomrule
\end{tabularx}
\end{center}

\begin{ideabox}
Ce consensus n'est pas anodin. Il traverse les clivages politiques, les frontières disciplinaires et les intérêts industriels.
\end{ideabox}

\clearpage

% ============================================================
% PARTIE II
% ============================================================
\part{Notre mission, notre vision, nos valeurs}
\label{part:mission}

\chapter{Mission et vision}

\section{Mission}

\begin{ctabox}
\begin{center}
{\Large\textbf{\faBullseye\ \ Notre mission}}\\[8pt]
{\large Rendre l'IA \textbf{équitable, transparente et responsable pour tous}, en fournissant les outils, les référentiels et les connaissances nécessaires à l'identification, la documentation et la correction des biais algorithmiques.}
\end{center}
\end{ctabox}

\section{Vision 2030}

\begin{ctabox}
\begin{center}
{\Large\textbf{\faBinoculars\ \ Notre vision}}\\[8pt]
{\large Un monde où \textbf{aucun système d'IA ne peut être déployé dans un secteur critique sans certification d'équité algorithmique}, de la même façon qu'aucun médicament ne peut être mis sur le marché sans essais cliniques rigoureux.}
\end{center}
\end{ctabox}

\section{Valeurs fondatrices}

\begin{center}
\renewcommand{\arraystretch}{1.4}
\begin{tabularx}{\textwidth}{@{}p{0.8cm}>{\bfseries}p{3.5cm}X@{}}
\toprule
\rowcolor{navy}
\textcolor{white}{} & \textcolor{white}{Valeur} & \textcolor{white}{Engagement} \\
\midrule
\faBalanceScale & Équité algorithmique & Aucun groupe ne doit être systématiquement désavantagé \\
\rowcolor{light}
\faSearch & Transparence radicale & Les mécanismes de décision doivent être explicables \\
\faShieldAlt & Indépendance & Aucune dépendance envers les acteurs que nous évaluons \\
\rowcolor{light}
\faGlobe & Inclusivité & Voix du Sud Global, des minorités, des communautés marginalisées \\
\faMicroscope & Rigueur scientifique & Toute affirmation est étayée par des preuves vérifiables \\
\bottomrule
\end{tabularx}
\end{center}

\clearpage

% ============================================================
% PARTIE III
% ============================================================
\part{Notre réponse}
\label{part:base}

\chapter{La Base IBI}

\section{Qu'est-ce que la Base IBI ?}

La \textbf{Base IBI} est la première base de données structurée au monde dédiée à la \textbf{cartographie exhaustive des biais dans l'intelligence artificielle}. Elle couvre :
\begin{itemize}[leftmargin=1.5em]
  \item les \textbf{types de biais} (taxonomie à 9 familles, 205+ types)
  \item les \textbf{secteurs impactés} (santé, justice, finance, RH, éducation, défense…)
  \item les \textbf{cas documentés} (50+ incidents réels avec analyse forensique)
  \item les \textbf{réglementations applicables} (63 cadres juridiques)
  \item les \textbf{mécanismes de mitigation} (80+ techniques et outils)
\end{itemize}

\section{Chiffres clés}

\begin{center}
\renewcommand{\arraystretch}{1.4}
\begin{tabularx}{\textwidth}{@{}>{\bfseries}p{5cm}>{\Large}cX@{}}
\toprule
\rowcolor{navy}
\textcolor{white}{Indicateur} & \textcolor{white}{Valeur} & \textcolor{white}{Détail} \\
\midrule
Biais documentés & \textbf{205+} & Types et sous-types catégorisés et décrits \\
\rowcolor{light}
Réglementations couvertes & \textbf{63} & EU AI Act, RGPD, NIS2, NIST AI RMF, ISO/IEC 42001… \\
Familles de biais & \textbf{9} & Taxonomie originale IBI \\
\rowcolor{light}
Secteurs couverts & \textbf{12+} & Santé, Justice, Finance, RH, Éducation, Défense… \\
Cas documentés & \textbf{50+} & Incidents réels avec analyse cause-effet-impact \\
\rowcolor{light}
Mécanismes de mitigation & \textbf{80+} & Techniques, procédures et outils de correction \\
\bottomrule
\end{tabularx}
\end{center}

\section{Les 9 familles de biais}

\begin{enumerate}[leftmargin=1.5em]
  \item \textbf{Biais de données} — sous-représentation ou sur-représentation dans les données d'entraînement
  \item \textbf{Biais d'annotation} — préjugés humains lors de l'étiquetage
  \item \textbf{Biais d'agrégation} — modèle entraîné sur une population hétérogène
  \item \textbf{Biais historique} — reproduction de discriminations passées
  \item \textbf{Biais de mesure} — proxys inadéquats (ex. code postal pour revenu)
  \item \textbf{Biais d'évaluation} — métriques masquant les disparités par sous-groupe
  \item \textbf{Biais de déploiement} — écart entre population d'entraînement et population réelle
  \item \textbf{Biais d'interaction} — comportements utilisateurs (RLHF, feedback)
  \item \textbf{Biais réglementaire} — lacunes normatives permettant des discriminations implicites
\end{enumerate}

\section{Accessibilité}

La Base IBI sera accessible via :

\begin{center}
\renewcommand{\arraystretch}{1.3}
\begin{tabularx}{\textwidth}{@{}>{\bfseries}p{3.5cm}X@{}}
\toprule
\rowcolor{navy}
\textcolor{white}{Canal} & \textcolor{white}{Description} \\
\midrule
\textbf{API REST v1.0} & Accès programmatique gratuit pour les ONG, la société civile et les chercheurs \\
\rowcolor{light}
\textbf{Interface web} & Recherche, filtres, export PDF, accessible sans compétence technique \\
\textbf{Intégration AEXIA\textsuperscript{TM}} & Connexion native avec les niveaux de maturité AEXIA-O et AEXIA-S \\
\bottomrule
\end{tabularx}
\end{center}

\clearpage

% ============================================================
% PARTIE IV
% ============================================================
\part{Gouvernance et engagements}
\label{part:gouvernance}

\chapter{Notre gouvernance}

\begin{center}
\renewcommand{\arraystretch}{1.3}
\begin{tabularx}{\textwidth}{@{}>{\bfseries}p{4.5cm}X@{}}
\toprule
\rowcolor{navy}
\textcolor{white}{Instance} & \textcolor{white}{Rôle} \\
\midrule
\textbf{Assemblée générale} & Décisions stratégiques, élection du Conseil d'administration \\
\rowcolor{light}
\textbf{Conseil d'administration} & Pilotage opérationnel, 5 à 9 membres élus \\
\textbf{Conseil scientifique} & Validation méthodologique, 7 à 12 experts internationaux \\
\rowcolor{light}
\textbf{Comité d'éthique} & Veille éthique, saisine citoyenne, 5 à 7 membres indépendants \\
\bottomrule
\end{tabularx}
\end{center}

\vspace{0.5cm}

\textbf{Principes de gouvernance} :
\begin{itemize}[leftmargin=1.5em]
  \item Indépendance totale vis-à-vis des acteurs évalués
  \item Publication annuelle des comptes et des financements
  \item Renouvellement régulier des instances
  \item Représentation équilibrée Nord/Sud, hommes/femmes, disciplines
\end{itemize}

\chapter{Rejoindre l'initiative}

\section{Quatre publics, quatre contributions}

IA BIAIS INITIATIVE s'adresse à \textbf{quatre publics complémentaires} :

\begin{center}
\renewcommand{\arraystretch}{1.4}
\small
\begin{tabularx}{\textwidth}{@{}>{\bfseries}p{4cm}X X@{}}
\toprule
\rowcolor{navy}
\textcolor{white}{Public} & \textcolor{white}{Contribution} & \textcolor{white}{Bénéfice} \\
\midrule
\textbf{Chercheurs \& académiques} & Documentation, publications, revue par les pairs & Co-autoriat, visibilité, accès API prioritaire \\
\rowcolor{light}
\textbf{Entreprises \& praticiens} \newline \footnotesize (CDO, DPO, DSI, Data Scientists) & Retours terrain, cas d'usage réels, audits & Outil d'audit gratuit, certification partenaire \\
\textbf{ONG \& société civile} & Collecte de cas, plaidoyer, sensibilisation & Ressources gratuites, plateforme de visibilité \\
\rowcolor{light}
\textbf{Régulateurs \& institutions} & Validation réglementaire, orientation stratégique & Influence sur les standards, consultation directe \\
\bottomrule
\end{tabularx}
\end{center}

\section{Comment nous rejoindre}

\begin{ctabox}
\begin{center}
{\Large\textbf{3 étapes pour rejoindre l'initiative}}\\[10pt]
\begin{tabular}{@{}c@{}}
\textbf{1. Manifestez votre intérêt} \quad $\rightarrow$ \quad contact@ia-biais-initiative.org\\[8pt]
\textbf{2. Signez la Charte d'engagement} \quad $\rightarrow$ \quad 2 pages\\[8pt]
\textbf{3. Contribuez activement} \quad $\rightarrow$ \quad cas, données, expertise, relais
\end{tabular}
\end{center}
\end{ctabox}

\vspace{0.5cm}

\section{Feuille de route 2026-2028}

\begin{center}
\renewcommand{\arraystretch}{1.3}
\small
\begin{tabularx}{\textwidth}{@{}>{\bfseries}p{3cm}p{2cm}X@{}}
\toprule
\rowcolor{navy}
\textcolor{white}{Phase} & \textcolor{white}{Période} & \textcolor{white}{Jalons clés} \\
\midrule
\textbf{Phase 0}\newline Fondation & T4 2026 & White Paper, création association, lancement site, connexion AEXIA\textsuperscript{TM} \\
\rowcolor{light}
\textbf{Phase 1}\newline Lancement & T1-T2 2027 & API Base IBI (beta), 3 partenariats académiques, 1 000 abonnés \\
\textbf{Phase 2}\newline Développement & T3-T4 2027 & 1er rapport annuel, formation certifiante, 5 000 utilisateurs API \\
\rowcolor{light}
\textbf{Phase 3}\newline Institutionnalisation & 2028 & Statut ONG, partenariats UE/UA, 1er Bias Summit Paris \\
\bottomrule
\end{tabularx}
\end{center}

\chapter{Conclusion}

Les leaders mondiaux de l'IA — de Sam Altman à Geoffrey Hinton, de Dario Amodei à Yoshua Bengio — s'accordent sur un point : \textbf{l'IA non maîtrisée représente un risque majeur pour l'humanité}. Les biais algorithmiques sont l'une des manifestations les plus concrètes, les plus immédiates et les plus injustes de ce risque.

\vspace{0.3cm}

IA BIAIS INITIATIVE ne prétend pas résoudre seule ce problème planétaire. Elle s'engage à être \textbf{la ressource de référence}, \textbf{l'espace de convergence} et \textbf{le catalyseur d'action} pour toutes celles et ceux qui refusent que l'IA reproduise — et amplifie — les injustices du monde.

\vspace{0.8cm}

\begin{ctabox}
\begin{center}
{\Large\itshape « L'IA est l'outil le plus puissant jamais créé par l'humanité.\\
La question n'est pas de savoir si nous allons l'utiliser —\\
elle est de savoir si nous allons le faire avec équité, transparence et responsabilité. »}\\[10pt]
{\large — \textbf{SIDIBE Mohamadou}, Fondateur}
\end{center}
\end{ctabox}

\clearpage

% ------------------------------------------------------------
% PAGE FINALE
% ------------------------------------------------------------
\backmatter

\chapter*{À propos du fondateur}
\addcontentsline{toc}{chapter}{À propos du fondateur}

\begin{center}

\begin{tikzpicture}
  \fill[purple] (0,0) rectangle (3,0.4);
  \fill[cyan]   (3,0) rectangle (6,0.4);
  \fill[orange] (6,0) rectangle (9,0.4);
  \fill[red]    (9,0) rectangle (12,0.4);
  \fill[green]  (12,0) rectangle (15,0.4);
\end{tikzpicture}

\vspace{0.5cm}

{\Large\textbf{Mohamadou SIDIBE}}

\vspace{0.3cm}

{\normalsize Expert en Gouvernance des Données \& IA}\\
{\normalsize Chief Data Officer · Polytechnique ExED 2026}\\[4pt]
{\normalsize\textbf{Fondateur, IA BIAIS INITIATIVE}}\\[2pt]
{\normalsize Auteur du framework AEXIA\textsuperscript{TM} v3.0}

\vspace{0.5cm}

{\small\texttt{contact@ia-biais-initiative.org}}\\[2pt]
{\small\href{https://ia-biais-initiative.org}{ia-biais-initiative.org}}

\vspace{0.5cm}

\begin{tikzpicture}
  \fill[purple] (0,0) rectangle (3,0.4);
  \fill[cyan]   (3,0) rectangle (6,0.4);
  \fill[orange] (6,0) rectangle (9,0.4);
  \fill[red]    (9,0) rectangle (12,0.4);
  \fill[green]  (12,0) rectangle (15,0.4);
\end{tikzpicture}

\end{center}

\vspace{1cm}

\begin{center}
{\small\textcolor{gray}{
  \textbf{IA BIAIS INITIATIVE} · Version 1.0 · Octobre 2026\\
  Licence CC BY 4.0 · Réutilisation libre avec attribution\\
  DOI du framework AEXIA\textsuperscript{TM} : \href{https://doi.org/10.5281/zenodo.23126471}{10.5281/zenodo.23126471}
}}
\end{center}

\end{document}